\scanpage{66}
\diagram{\deleted{\begin{tikzcd}[column sep=small]
V \arrow[r] & V_{0,1} & X_0^* \arrow[l]
\end{tikzcd}}}

\margin{\uline{NB} ou mieux
\diagram{\begin{tikzcd}[row sep=small,column sep=small]
X_0^* \arrow[d] \arrow[r,dash] & {} \\
V_{01} & V_{01}^* \arrow[l] \arrow[d] \\
{} & X_1^* & {} \arrow[l]
\end{tikzcd}}}

\diagram{(20)\qquad
\begin{tikzcd}[column sep=.3em,row sep=1.2em]
X_0^* \arrow[r] & \underset{\deleted{\downarrow}}{V_{0,1}^{\unsure{*|\deleted{*}}}} & V_{0,1}^* \arrow[l] \arrow[d] & & & & & & & \\
& & X_1^* \arrow[r] & V_{1,2} & V_{1,2}^* \arrow[l] \arrow[d] & & & & & \\
& & & & X_2^* \arrow[r] & V_{2,3} & V_{2,3}^* \arrow[l] \arrow[d] & & & \\
& & & & & & X_3^* & & & \\
& & & \widetilde\Delta_n & & & \mathord{-}\;\mathord{-}\;\mathord{-}\;\mathord{-} & {} \arrow[d] & & \\
& & & & & & & X_{n-1}^* \arrow[r] & V_{n-1,n} & V_{n-1,n}^* \arrow[l] \arrow[d] \\
& & & & & & & & & X_n^*.
\end{tikzcd}}

Si $I$ est le diagramme
\diagram{\begin{tikzcd}[column sep=small,row sep=small]
& a \\
\unsure{\circ|\bullet} \arrow[ru] \arrow[rd] & \\
& b
\end{tikzcd}}, on trouve
\diagram{(21)\qquad
\begin{tikzcd}[column sep=small,row sep=large]
\deleted{\unsure{X_b^*|X_a^*}} & V_{0,b}^* \arrow[r] \arrow[d] & V_{0,b} & X_0^* \arrow[l] \arrow[r] & V_{0,a} & V_{0,a}^* \arrow[l] \arrow[d] \\
& X_b^* & & & & X_a^*.
\end{tikzcd}}

Si $I$ est le diagramme
\diagram{\begin{tikzcd}[column sep=small,row sep=small]
a\mathord\cdot \arrow[rd] & \\
& c \\
b\mathord\cdot \arrow[ru,leftrightarrow] &
\end{tikzcd}}, on trouve
\diagram{(22)\qquad
\begin{tikzcd}[column sep=large,row sep=large]
X_a^* \arrow[d] & & X_b^* \arrow[d] \\
V_{a,c} \arrow[d,leftrightarrow] & & V_{b,c} \\
V_{a,c}^* \arrow[r] & X_c^* & V_{b,c}^* \arrow[u] \arrow[l]
\end{tikzcd}},
etc.

\deleted{La première chose qu'on aimerait préciser est le} Supposons qu'on dispose sur $I$ d'une fonction à valeurs entières

\margin{65}
